\section{Bundles, tangent bundles and differential forms}\label{sec:geometry}

\subsection{From a bad operator family to changes of frame}

Transition functions take values in invertible operators; for a single Banach fiber, invertible families impose the same condition as all operator families (Proposition~\ref{fam2:prop:reductions}(1)).

\begin{proposition}[shear realization]\label{prop:shear}
Suppose $A=A^k_{E,D;F}$ is not analytic at $u_0\in\cL(E,D)$. Set $H=D\oplus E$, with the maximum norm, and
\[
 g(u)(d,e)=(d+ue,e).
\]
Both $g$ and its inverse family are affine analytic. Nevertheless, the family of pullback operators $g(u)^*$ on $\Alt^k(H;F)$ is not analytic at $u_0$.
\end{proposition}

\begin{proof}
The inverse family is $g(-u)$. Let $\pi_D\colon H\to D$ be the projection and $i_E(e)=(0,e)$. Define bounded linear maps $J\colon\Alt^k(D;F)\to\Alt^k(H;F)$ by $Jm=m\circ\pi_D^{\times k}$ and $R\colon\Alt^k(H;F)\to\Alt^k(E;F)$ by restriction along $i_E$. Since $\pi_Dg(u)i_E=u$,
\begin{equation}\label{eq:shear-compression}
 R\circ g(u)^*\circ J=A(u).
\end{equation}
Bounded linear compression preserves analyticity, so $u\mapsto g(u)^*$ is not analytic at $u_0$.
\end{proof}

\begin{corollary}[shear bundle]\label{cor:shear-bundle}
Suppose $A=A^k_{E,D;F}$ is not analytic at $u_0\in\cL(E,D)$, and let $H$ and $g$ be as in Proposition~\ref{prop:shear}. Take the base $M=\cL(E,D)$ and the trivial bundle $M\times H$. Its product trivialization and the trivialization $(u,z)\mapsto(u,g(u)z)$ form an analytic bundle atlas, and the induced alternating atlas on the trivial bundle with fiber $\Alt^k(H;F)$ is not analytic. Equivalently, the analytic bundle automorphism $(u,z)\mapsto(u,g(u)z)$ does not induce an analytic morphism on the alternating product bundle.
\end{corollary}

\begin{proof}
The transitions are $g(u)$ and $g(-u)$, affine analytic by Proposition~\ref{prop:shear}. The induced alternating transitions are pullback by $g(u)$ and $g(-u)$, and the induced morphism is locally pullback by $g(u)$; by Proposition~\ref{prop:shear}, $u\mapsto g(u)^*$ is not analytic at $u_0$.
\end{proof}

Thus every failure of the hom-space criterion has a bundle realization, and Theorem~\ref{thm:operator-obstructions} gives fibers for which it occurs at every point of the base. Applied to analytic families $x\mapsto u(x)$, the shear also makes the familywise hypothesis of Theorem~\ref{thm:bundle-lifting} necessary for all fibers and frames; Section~\ref{sec:families} studies it (Corollary~\ref{fam2:cor:bundles}, Theorems~\ref{fam2:thm:polarizable} and~\ref{fam2:thm:bases}).

\subsection{Tangent bundles and individual forms}

\begin{proposition}[tangent and pullback obstructions]\label{prop:tangent-obstruction}
Suppose $K$ is complete, has characteristic $p>0$, is not spherically complete, and $k\ge p$.
\begin{enumerate}
\item There is an analytic Banach manifold with two global analytic charts whose induced transition on scalar alternating $k$\nobreakdash-forms is analytic at no point.
\item There are Banach spaces $N,Y$, a polynomial analytic map $h\colon N\to Y$, and an analytic map $\omega_0\colon Y\to\Alt^k(Y;K)$ such that $h^*\omega_0$, expressed in the product coordinates on $N$, is analytic at no point as an $\Alt^k(N;K)$\nobreakdash-valued map. After inclusion in $\Mult^k(N;K)$ it is analytic everywhere.
\end{enumerate}
\end{proposition}

The proof is in Appendix~\ref{app:deferred}.

Assertion~(2) separates intrinsic analyticity of individual forms from analyticity of bundle transitions; Section~\ref{sec:forms} builds a calculus on the ambient notion.

The chart $\psi$ of the proof of Proposition~\ref{prop:tangent-obstruction} realizes any nowhere analytic action $A^k_{E,D;F}$; with Theorem~\ref{thm:operator-obstructions}(1) it works over every field of characteristic $p$.
\begin{proposition}[tangent obstructions with other values]\label{prop:tangent-values}
Let $K$ be a nontrivially normed field of characteristic $p>0$, complete or not, and let $k\ge p$.
\begin{enumerate}
\item There are a Banach space $F$ and an analytic Banach manifold with two global analytic charts whose induced transition on $F$\nobreakdash-valued alternating $k$\nobreakdash-forms is analytic at no point.
\item There is an analytic Banach manifold with two global analytic charts whose induced transition on tangent-valued alternating $k$\nobreakdash-forms is analytic at no point.
\item If $K$ is complete and not spherically complete, a manifold as in~(2) exists with a nonarchimedean Banach model.
\end{enumerate}
The target $F$ constructed for~(1) admits no equivalent nonarchimedean norm; for $K=\F_q((X))$ as in Corollary~\ref{lau:cor:dichotomy}, the model space can be chosen nonarchimedean. The model space constructed for~(2) admits no equivalent nonarchimedean norm.
\end{proposition}
The proof is in Appendix~\ref{app:deferred}.

\paragraph{\textit{Bourbaki's question.}}
For $k\ge1$, $\Alt^k$ on Banach spaces is the functor $\alpha_k$ of~\cite[n\textsuperscript{o}~7.8.1]{Bou}, and $A^k_{E,D;F}(u)=\alpha_k(\id_F,u)$. There $K$ is $\mathbb R$, $\mathbb C$ or complete nonarchimedean, norms over a nonarchimedean $K$ are nonarchimedean, and class $C^r$ means analytic if $K\ne\mathbb R$, in a sense that agrees with ours between Banach spaces. In that setting $\alpha_k$ is a vector functor of class $C^r$ if $k!\ne0$ in $K$ (Corollary~\ref{cor:classification}), and for every $k$ if $K$ is discretely valued (Corollary~\ref{cor:discrete-targets} and the proof of Corollary~\ref{cor:bundle-cases}). It is not one if $k!=0$ in $K$ and $K$ is not spherically complete, since Theorem~\ref{thm:operator-obstructions}(2) gives nonarchimedean Banach spaces $E,D$ with $A^k_{E,D;K}$ analytic at no point; so the definition in~\cite[n\textsuperscript{o}~7.8.1]{Bou} needs a restriction. By Corollary~\ref{cor:shear-bundle} and Proposition~\ref{prop:tangent-obstruction}, it is needed for the bundles themselves. Source spaces of countable type suffice in every degree, by the following corollary of Theorem~\ref{fib:thm:ultra}.

\begin{corollary}[countable type and Bourbaki's $\alpha_k$]\label{fib:cor:bourbaki}
Let $K$ be complete and nonarchimedean, let $E,D$ be normed spaces, at least one of them of countable type, and let
$F$ be a nonarchimedean Banach space. For every normed $F'$, the joint action $B$ of
Proposition~\ref{prop:ambient-polynomial} on $\cL(E,D)\times\cL(F,F')$ has a bounded $(k+1)$\nobreakdash-linear lift.
In the setting of the paragraph \emph{Bourbaki's question}, $\alpha_k$ is therefore a
vector functor of the class considered there on source spaces of countable type, in every degree.
\end{corollary}

\begin{proof}
If $Q$ lifts $A^k_{E,D;F}$ (Theorem~\ref{fib:thm:ultra}), then $(u_1,\dots,u_k,g)\mapsto(m\mapsto g\circ Q(u_1,\dots,u_k)(m))$
lifts $B$, as in the proof of Corollary~\ref{cor:bundle-cases}.
\end{proof}

Our results do not decide the case of a densely valued, spherically complete field with $k!=0$ (Section~\ref{sec:questions}). Theorem~\ref{amb:thm:seven-eight} gives the constructions of~\cite[n\textsuperscript{o}~7.8.1--7.8.4]{Bou} on ambient sections of filtered analytic bundles, in every characteristic.

\paragraph{\textit{A derivative adapted to positive characteristic.}}
Ya\"el Dillies asked whether $A=A^k_{E,D;F}$ is analytic in the calculus of~\cite[Definition~4.7 and Theorem~5.4]{BGN}, where a $C^k$ map $f$ has Taylor coefficients $a_n(x,h)$ with $n!\,a_n(x,h)=d^nf(x)(h,\ldots,h)$. There $A$ is smooth~\cite[Remark~4.2 and Lemma~10.1]{BGN} and $a_k(u_0,u)=A(u)$, while $d^kA(u_0)(u,\ldots,u)=k!\,A(u)$ vanishes if $k!=0$. Analyticity is not restored: it asks for a bounded $k$-linear lift (Proposition~\ref{prop:homogeneous-lift}), and Theorem~\ref{thm:operator-obstructions} gives spaces with none when $k!=0$ in $K$. Section~\ref{amb:sec} enlarges the target instead; \cite{II} changes the coefficients.

\subsection{Analytic differential forms}\label{sec:forms}

Throughout this subsection $K$ is complete and nontrivially normed, manifold models are normed spaces, and charts and maps are $C^\omega$, so that their derivatives are analytic; on the Banach coefficient spaces this is power-series analyticity.

\begin{definition}
An \emph{ambient analytic scalar $k$\nobreakdash-form} on $M$ is a family $\eta_x\in\Alt^k(T_xM;K)$ whose representative in each chart is analytic after the inclusion
\[
 j\colon\Alt^k(P;K)\hookrightarrow\Mult^k(P;K).
\]
Write $\Omega^k_{\mathrm{amb}}(M)$ for the space of such forms. In degree zero this is the space of analytic scalar functions.
\end{definition}

Under a coordinate change $\psi=\phi'\circ\phi^{-1}$ the representatives satisfy $j\eta_\phi(y)=(j\eta_{\phi'}(\psi(y)))\circ(D\psi(y),\ldots,D\psi(y))$, which is analytic by bounded multilinear composition; so any analytic atlas may be used.

For an ambient $k$\nobreakdash-form $\eta$, an ambient $l$\nobreakdash-form $\zeta$ and $S\subseteq\{1,\ldots,k+l\}$, let $v_S$ list the arguments $v_i$, $i\in S$, in increasing order of $i$, and let $\epsilon(S)$ be the shuffle sign of Section~\ref{cmb:ssec:degrees}. The product is the shuffle formula
\begin{equation}\label{eq:form-shuffle}
 (\eta\wedge\zeta)(v_1,\ldots,v_{k+l})
 =\sum_{\abs{S}=k}\epsilon(S)\eta(v_S)\zeta(v_{S^c}),
\end{equation}
a bounded bilinear operation on multilinear coefficient spaces, of norm at most $\binom{k+l}{k}$, which is alternating on alternating inputs. On an open subset of a model space $P$, the exterior derivative of an ambient coefficient map $\eta$ of degree $k$ is
\begin{equation}\label{eq:exterior-derivative}
 (d\eta)_y(v_0,\ldots,v_k)
 =\sum_{i=0}^k(-1)^iD\eta(y)(v_i)
       (v_0,\ldots,\widehat v_i,\ldots,v_k).
\end{equation}

\begin{theorem}[ambient differential calculus]\label{thm:ambient-forms}
Ambient analytic forms define a sheaf of graded algebras on every analytic manifold. Analytic maps induce pullbacks, and there is an exterior derivative
\[
 d\colon\Omega^k_{\mathrm{amb}}(M)\longrightarrow\Omega^{k+1}_{\mathrm{amb}}(M).
\]
The product is associative and graded commutative, $d^2=0$, pullback commutes with both product and exterior derivative, and
\[
 d(\eta\wedge\zeta)=d\eta\wedge\zeta+(-1)^k\eta\wedge d\zeta
 \qquad(\eta\in\Omega^k_{\mathrm{amb}}(M)).
\]
Thus $\Omega^\bullet_{\mathrm{amb}}$ is a contravariant differential graded algebra construction in every characteristic.
\end{theorem}
The proof is in Appendix~\ref{app:deferred}.

\begin{proposition}[comparison with alternating-bundle sections]\label{prop:ambient-agreement}
Every intrinsically analytic alternating-valued coefficient map is ambient analytic. If, for the model space $P$, the inclusion $j$ has a bounded linear retraction, then the natural alternating cotangent atlas is analytic and its analytic sections are precisely the ambient analytic forms. The same conclusion holds when $P$ has finite continuous coordinates.
\end{proposition}

\begin{proof}
The first statement follows by composition with $j$. For a bounded retraction $R$, the identity $\eta=R(j\eta)$ gives the converse. Moreover, $(u_1,\ldots,u_k)\mapsto\bigl[m\mapsto R\bigl(m\circ(u_1,\ldots,u_k)\bigr)\bigr]$ is a bounded multilinear lift of the precomposition action, so the induced alternating cotangent transitions are analytic. If $P$ has finite continuous coordinates $(e_i,e^*_i)_{i\le d}$, the values $m(e_{i_1},\ldots,e_{i_k})$ identify $\Mult^k(P;K)$ with $K^{d^k}$, with its maximum norm, by a bounded linear isomorphism, so every linear retraction of $j$ is bounded and the previous case applies.
\end{proof}

The retraction exists if $k!$ is invertible, by normalized antisymmetrization, and if $K$ is spherically complete and nonarchimedean, by Theorem~\ref{thm:spherical}; and if $P$ has finite continuous coordinates, by the proof above.

\paragraph{\textit{Exact constant forms in degree $p$.}}
Here $K$ is any nontrivially normed field of characteristic $p>0$, $Y$ and $F$ are normed, and $\mathcal A_Y\colon\Mult^p(Y;F)\to\Alt^p(Y;F)$ is unnormalized antisymmetrization (Lemma~\ref{ob:lem:target}). Define $(d\eta)_y\in\Alt^p(Y;F)$ by~\eqref{eq:exterior-derivative} for every map $\eta$ from an open subset of $Y$ to $\Alt^{p-1}(Y;F)$ differentiable at $y$. Constant forms are closed, and $(\mathrm R^p_Y)$ is the exactness of one of them.
\begin{proposition}[exact constant forms]\label{fib:prop:derham}
For $\omega\in\Alt^p(Y;F)$ the following are equivalent:
\begin{enumerate}
\item there are a map $\eta$ from an open subset of $Y$ to $\Alt^{p-1}(Y;F)$ and a point $y$ at which $\eta$ is
differentiable with $(d\eta)_y=\omega$;
\item $\omega=d\eta$ on $Y$ for a bounded linear map $\eta\colon Y\to\Alt^{p-1}(Y;F)$;
\item $\omega=\mathcal A_Y\beta$ for some $\beta\in\Mult^p(Y;F)$.
\end{enumerate}
For $F=\Lambda^p_\pi Y$ and $\omega=\omega_Y$ they hold if and only if $\operatorname{alt}^p_Y$ has a bounded linear left
inverse, and then $(E,Y)$ is good in degree $p$ for every normed $E$. For $F=\widehat\Lambda^p_\pi Y$ and
$\omega=\omega_Y$ they hold if and only if $(\mathrm R^p_Y)$ holds, and then $(E,Y)$ is good for Banach values in degree
$p$ for every normed $E$.
\end{proposition}

The proof is in Appendix~\ref{app:deferred}.

Under condition~\textup{(L)} of Section~\ref{obs:ssec:field}, for $Y=c_0(\N,K)$ the closed constant form $\omega_Y$ with values in $\widehat\Lambda^p_\pi Y$ has no primitive near any point, by Theorem~\ref{fib:thm:ell2}.
